\begin{theorem}[Polynomial accuracy at a prescribed exponential rate]%
    \label{thm:ldp_polynomial_accuracy_of_rate}
    \lean{MPSTensor.exists_approximationError_le_polynomial_of_rate}
    \leanok
    Let $A$ be a normal tensor in the gauge $\sum_i(A^i)^\dagger A^i=\Id$,
    with a positive definite fixed matrix $\rho$ satisfying
    $\E_A(\rho)=\rho$ and $\Tr\rho=1$.
    Let $\lambda_2$ bound the moduli of the transfer eigenvalues other than
    $1$, assume $\xi=-1/\log|\lambda_2|>0$, and fix $0<\gamma<1/2$.
    There is a constant $C_\gamma>0$, independent of $\eta,q,M$, such
    that for every $\eta>0$ and every uniform decomposition $N=Mq\geq2$,
    \begin{align}
        q & =\left\lceil\frac{\xi(1+\eta)\log N}{\gamma}\right\rceil
        \quad\Longrightarrow\quad
        \varepsilon(\widetilde\phi_N,\phi_N)\leq C_\gamma N^{-\eta}.
        \label{eq:ldp_polynomial_prescribed_rate}
    \end{align}
    This is a polynomial-accuracy corollary of the exponential approximation
    estimate of~\cite{Malz2024Preparation}. The uniform-blocking and
    positive-correlation-length scope is recorded
    in~\cite{gap:mswc24_polynomial_accuracy_uniform_blocks}.
\end{theorem}
\begin{proof}\leanok
    \uses{thm:ldp_approximation_error_normal}
    Apply the second-order estimate at rate $\gamma/2$ to obtain
    $\varepsilon\leq C_\gamma M e^{-\gamma q/\xi}$.
    The ceiling condition gives $\gamma q\geq\xi(1+\eta)\log N$.
    Since $q\geq1$ and $M\leq N$,
    $C_\gamma M e^{-\gamma q/\xi}\leq C_\gamma N e^{-(1+\eta)\log N}
        =C_\gamma N^{-\eta}$.
\end{proof}

\begin{theorem}[Polynomial approximation accuracy under uniform blocking]%
    \label{thm:ldp_polynomial_accuracy_uniform_blocks}
    \lean{MPSTensor.exists_approximationError_le_polynomial_of_uniformBlocks}
    \leanok
    Let $A$ be a normal tensor in the gauge $\sum_i(A^i)^\dagger A^i=\Id$,
    with a positive definite fixed matrix $\rho$ satisfying
    $\E_A(\rho)=\rho$ and $\Tr\rho=1$.
    Let $\lambda_2$ bound the moduli of the transfer eigenvalues other than
    $1$, and assume $\xi=-1/\log|\lambda_2|>0$.
    There is a constant $C>0$ such that for every $\eta>0$ and every
    uniform decomposition $N=Mq\geq2$ with
    \begin{align}
        q & =\left\lceil2\xi(1+\eta)\log N\right\rceil,
        \label{eq:ldp_polynomial_block_length}
    \end{align}
    the fixed-point approximation satisfies
    \begin{align}
        \varepsilon(\widetilde\phi_N,\phi_N) & \leq C N^{-\eta}.
        \label{eq:ldp_polynomial_uniform_accuracy}
    \end{align}
    The constant is independent of $\eta,q,M$.
    This is the uniform-blocking, positive-correlation-length case of
    the paragraph following Lemma~1 of~\cite{Malz2024Preparation};
    the other chain lengths and zero correlation length are discussed
    in~\cite{gap:mswc24_polynomial_accuracy_uniform_blocks}.
\end{theorem}
\begin{proof}\leanok
    \uses{thm:ldp_approximation_error_normal}
    Apply the second-order approximation estimate at the fixed rate
    $\gamma=1/4$, obtaining $\varepsilon\leq C M e^{-q/(2\xi)}$.
    Formula~(\ref{eq:ldp_polynomial_block_length}) gives
    $q\geq2\xi(1+\eta)\log N$. Since $N=Mq\geq2$, one has
    $q\geq1$ and $M\leq N$. Thus
    \begin{align}
        C M e^{-q/(2\xi)}
         & \leq C N e^{-(1+\eta)\log N}=C N^{-\eta}.
        \notag
    \end{align}
    This proves~(\ref{eq:ldp_polynomial_uniform_accuracy}).
    The second-order rate is the justification for the exact coefficient
    in~(\ref{eq:ldp_polynomial_block_length}).
\end{proof}
